% Part: first-order-logic
% Chapter: natural-deduction
% Section: provability-quantifiers

% verification of properties of provability needed for maximally
% consistent sets in the completeness chapter.

\documentclass[../../../include/open-logic-section]{subfiles}

\begin{document}

\olfileid{fol}{ntd}{qpr}

\olsection{\usetoken{S}{derivability} en de kwantoren}

\begin{explain}
  Voor de volledigheidsstelling is tevens vereist dat de regels voor
  natuurlijke deductie de in deze paragraaf vastgestelde feiten over
  ~$\Proves$ opleveren.
\end{explain}

\begin{thm}
\ollabel{thm:strong-generalization} Als $c$ een constante is die niet
in $\Gamma$ of $!A(x)$ voorkomt en $\Gamma \Proves !A(c)$, dan $\Gamma
\Proves \lforall[x][!A(x)]$.
\end{thm}

\begin{proof}
Laat $\delta$ !!a{derivation} van $!A(c)$ uit $\Gamma$ zijn. Door een
\Intro{\lforall}-inferentie toe te voegen, verkrijgen we
!!a{derivation} van $\lforall[x][!A(x)]$. Omdat $c$ niet in $\Gamma$
of $!A(x)$ voorkomt, is aan de eigenvariabelevoorwaarde voldaan.
\end{proof}

\begin{prop}
\ollabel{prop:provability-quantifiers}
\begin{tagenumerate}{prvEx,prvAll}
\tagitem{prvEx}{$!A(t) \Proves \lexists[x][!A(x)]$.}{}

\tagitem{prvAll}{$\lforall[x][!A(x)] \Proves !A(t)$.}{}
\end{tagenumerate}
\end{prop}

\begin{proof}
\begin{tagenumerate}{prvEx,prvAll}
\tagitem{prvEx}{Het volgende is !!a{derivation}
  van~$\lexists[x][!A(x)]$ uit~$!A(t)$:
\begin{prooftree}
\AxiomC{$!A(t)$}
\RightLabel{\Intro{\lexists}}
\UnaryInfC{$\lexists[x][!A(x)]$}
\end{prooftree}}{}

\tagitem{prvAll}{Het volgende is !!a{derivation} van~$!A(t)$
  uit~$\lforall[x][!A(x)]$:
\begin{prooftree}
\AxiomC{$\lforall[x][!A(x)]$}
\RightLabel{\Elim{\lforall}}
\UnaryInfC{$!A(t)$}
\end{prooftree}}{}
\end{tagenumerate}
\end{proof}
\end{document}
